\documentclass[tikz,border=8pt]{standalone}
\usepackage{tikz}
\usepackage{amsmath}
\usepackage{amssymb}
\usetikzlibrary{
    positioning,
    shapes.geometric,
    shapes.multipart,
    arrows.meta,
    calc,
    fit,
    backgrounds,
    decorations.pathreplacing,
    patterns,
    shadows.blur,
    matrix,
    chains,
    fadings
}

% Publication color palette (colorblind-safe, Nature Methods style)
\definecolor{recv}{RGB}{230, 159, 0}       % Orange - receiver
\definecolor{sender}{RGB}{0, 158, 115}     % Teal - senders/niche
\definecolor{hlca}{RGB}{86, 180, 233}      % Sky blue - HLCA reference
\definecolor{luca}{RGB}{213, 94, 0}        % Vermillion - LuCA reference
\definecolor{lr}{RGB}{204, 121, 167}       % Pink - L-R tokens
\definecolor{relay}{RGB}{117, 112, 179}    % Purple - relay tokens
\definecolor{attn}{RGB}{240, 228, 66}      % Yellow - attention
\definecolor{output}{RGB}{0, 114, 178}     % Blue - outputs
\definecolor{bglight}{RGB}{250, 250, 252}  % Light background
\definecolor{bggray}{RGB}{240, 240, 245}   % Gray background

\begin{document}
\begin{tikzpicture}[
    % === NODE STYLES ===
    token/.style={
        rectangle, draw=#1!80!black, fill=#1!40,
        rounded corners=2pt, minimum width=0.85cm, minimum height=0.6cm,
        font=\tiny\sffamily\bfseries, line width=0.6pt, align=center
    },
    smalltoken/.style={
        rectangle, draw=#1!80!black, fill=#1!30,
        rounded corners=1.5pt, minimum width=0.55cm, minimum height=0.45cm,
        font=\tiny\sffamily, line width=0.5pt
    },
    encoder/.style={
        rectangle, draw=#1!70!black, fill=#1!15,
        rounded corners=4pt, minimum width=2.8cm, minimum height=1.1cm,
        font=\scriptsize\sffamily\bfseries, line width=1pt, align=center,
        blur shadow={shadow blur steps=5, shadow xshift=0.4pt, shadow yshift=-0.4pt}
    },
    attnblock/.style={
        rectangle, draw=attn!80!black, fill=attn!20,
        rounded corners=3pt, minimum width=2.4cm, minimum height=0.7cm,
        font=\tiny\sffamily\bfseries, line width=0.8pt, align=center
    },
    ffn/.style={
        rectangle, draw=gray!60, fill=gray!10,
        rounded corners=2pt, minimum width=2.4cm, minimum height=0.5cm,
        font=\tiny\sffamily, line width=0.6pt
    },
    layernorm/.style={
        rectangle, draw=gray!50, fill=white,
        rounded corners=1pt, minimum width=0.8cm, minimum height=0.3cm,
        font=\tiny\sffamily, line width=0.4pt
    },
    projection/.style={
        trapezium, trapezium angle=70, draw=#1!70!black, fill=#1!25,
        minimum width=1.2cm, minimum height=0.5cm,
        font=\tiny\sffamily, line width=0.6pt
    },
    head/.style={
        rectangle, draw=output!80!black, fill=output!25,
        rounded corners=3pt, minimum width=1.5cm, minimum height=0.6cm,
        font=\tiny\sffamily\bfseries, line width=0.7pt
    },
    annotbox/.style={
        rectangle, draw=gray!40, fill=bglight,
        rounded corners=2pt, font=\tiny\sffamily, align=left,
        text width=#1, inner sep=3pt
    },
    arr/.style={-{Stealth[length=2mm, width=1.2mm]}, line width=0.6pt, #1},
    arrdash/.style={-{Stealth[length=1.5mm, width=1mm]}, line width=0.5pt, dashed, #1},
    brace/.style={decorate, decoration={brace, amplitude=3pt, raise=1pt}},
    sectitle/.style={font=\scriptsize\sffamily\bfseries, text=gray!70!black}
]

% ============================================
% TITLE
% ============================================
\node[font=\normalsize\sffamily\bfseries] at (5, 13.5) {StageBridge Transformer Architecture};
\node[font=\tiny\sffamily, text=gray!60] at (5, 13.1) {Receiver-Centered Communication Encoder with L-R Integration};

% ============================================
% SECTION 1: INPUT TOKENS (top)
% ============================================
\node[sectitle] at (-0.5, 12.3) {Input Tokens};

% Background box for tokens
\begin{scope}[on background layer]
    \fill[bggray, rounded corners=5pt] (-0.8, 11.2) rectangle (10.8, 12.1);
\end{scope}

% Token sequence
\node[token=recv] (recv) at (0, 11.6) {Recv};
\node[token=sender] (s1) at (1.1, 11.6) {$S_1$};
\node[token=sender!80] (s2) at (2.0, 11.6) {$S_2$};
\node[token=sender!60] (s3) at (2.9, 11.6) {$\cdots$};
\node[token=sender!40] (sk) at (3.8, 11.6) {$S_k$};
\node[token=lr] (lr1) at (5.2, 11.6) {LR$_1$};
\node[token=lr!80] (lr2) at (6.1, 11.6) {LR$_2$};
\node[token=lr!60] (lrn) at (7.0, 11.6) {$\cdots$};
\node[token=relay] (rl1) at (8.2, 11.6) {Rly$_1$};
\node[token=relay!80] (rl2) at (9.1, 11.6) {Rly$_2$};
\node[token=relay!60] (rln) at (10.0, 11.6) {$\cdots$};

% Token group labels
\node[font=\tiny\sffamily, below=0.15cm of recv] {Receiver};
\node[font=\tiny\sffamily, below=0.15cm of s2] {Senders};
\node[font=\tiny\sffamily, below=0.15cm of lr2] {L-R Pairs};
\node[font=\tiny\sffamily, below=0.15cm of rl1] {Relay};

% Dimension annotations
\node[font=\tiny\ttfamily, text=recv!80!black, above=0.1cm of recv] {$d$};
\node[font=\tiny\ttfamily, text=sender!80!black, above=0.1cm of s2] {$k \times d$};
\node[font=\tiny\ttfamily, text=lr!80!black, above=0.1cm of lr2] {$n \times d$};
\node[font=\tiny\ttfamily, text=relay!80!black, above=0.1cm of rl1] {$m \times d$};

% ============================================
% SECTION 2: PROJECTION LAYERS
% ============================================
\node[sectitle] at (-0.5, 10.5) {Projections};

\node[projection=recv, minimum width=0.8cm] (proj_recv) at (0, 9.9) {$W_r$};
\node[projection=sender, minimum width=2.5cm] (proj_send) at (2.45, 9.9) {$W_s$ + Type + Ring + Offset};
\node[projection=lr, minimum width=1.8cm] (proj_lr) at (6.1, 9.9) {$W_{lr}$ + Type};
\node[projection=relay, minimum width=1.8cm] (proj_relay) at (9.1, 9.9) {$W_{rl}$ + Type};

% Arrows from tokens to projections
\draw[arr=recv!70] (recv) -- (proj_recv);
\draw[arr=sender!70] (s2) -- (proj_send);
\draw[arr=lr!70] (lr2) -- (proj_lr);
\draw[arr=relay!70] (rl1) -- (proj_relay);

% ============================================
% SECTION 3: SENDER ENCODER (Self-Attention)
% ============================================
\node[sectitle] at (-0.5, 8.8) {Sender Encoder};

\node[encoder=sender, minimum width=3.5cm] (sender_enc) at (2.45, 7.8) {
    Self-Attention Block\\[-2pt]
    \tiny SAB: Set Attention
};

\draw[arr=sender, line width=0.8pt] (proj_send) -- (sender_enc);

% Internal detail box
\node[annotbox=2cm, right=0.3cm of sender_enc] at (5.5, 7.8) {
    \textbf{Sender Self-Attn:}\\
    $H_s = \text{SAB}(S)$\\
    Encodes sender\\
    cell relationships
};

% ============================================
% SECTION 4: COMMUNICATION ENCODER (Cross-Attention)
% ============================================
\node[sectitle] at (-0.5, 6.5) {Communication Encoder};

\node[encoder=lr, minimum width=4.5cm] (comm_enc) at (4, 5.5) {
    Cross-Attention Block\\[-2pt]
    \tiny Query: LR tokens, Key/Value: Senders
};

\draw[arr=lr, line width=0.8pt] (proj_lr) -- ++(0, -2.5) -- (comm_enc.north-|6.1,0);
\draw[arr=sender, line width=0.8pt] (sender_enc) -- ++(0, -0.7) -| (comm_enc);

% Cross attention formula
\node[annotbox=3cm] at (9, 5.5) {
    \textbf{L-R $\times$ Sender:}\\
    $\text{Attn}_{lr\leftarrow s} = \text{softmax}\left(\frac{Q_{lr}K_s^\top}{\sqrt{d}}\right)V_s$\\[2pt]
    L-R tokens query sender\\
    memory for interaction
};

% ============================================
% SECTION 5: RELAY ENCODER (Self-Attention)
% ============================================
\node[sectitle] at (6.5, 4.2) {Relay Encoder};

\node[encoder=relay, minimum width=3cm] (relay_enc) at (9.1, 3.2) {
    Self-Attention Block\\[-2pt]
    \tiny Response + Relay tokens
};

\draw[arr=relay, line width=0.8pt] (proj_relay) -- ++(0, -5.7) -- (relay_enc.north);

% ============================================
% SECTION 6: MEMORY BANK AGGREGATION
% ============================================
\node[sectitle] at (-0.5, 3.8) {Memory Bank};

% Memory aggregation
\begin{scope}[on background layer]
    \fill[attn!10, rounded corners=4pt] (-0.5, 2.0) rectangle (7, 3.5);
\end{scope}

\node[smalltoken=sender] (mem_s) at (0.3, 2.7) {$H_s$};
\node[smalltoken=lr] (mem_lr) at (1.3, 2.7) {$H_{lr}$};
\node[smalltoken=relay] (mem_rl) at (2.3, 2.7) {$H_{rl}$};
\node[font=\scriptsize\sffamily] at (3.5, 2.7) {$\rightarrow$ Concat $\rightarrow$};
\node[token=attn, minimum width=1.5cm] (memory) at (5.5, 2.7) {Memory $M$};

\draw[arr=sender!60] (sender_enc) -- ++(0, -3.5) -| (mem_s);
\draw[arr=lr!60] (comm_enc) -- ++(0, -1.5) -| (mem_lr);
\draw[arr=relay!60] (relay_enc) -- ++(0, -0.3) -| (mem_rl.east);

\node[font=\tiny\sffamily, text=gray] at (3, 2.2) {Aggregate all encoded context};

% ============================================
% SECTION 7: RECEIVER QUERY (Cross-Attention)
% ============================================
\node[sectitle] at (-0.5, 1.3) {Receiver Query};

\node[encoder=recv, minimum width=5cm] (recv_query) at (3, 0.3) {
    Receiver $\leftarrow$ Memory Cross-Attention\\[-2pt]
    \tiny Query: Receiver, Key/Value: Memory Bank
};

\draw[arr=recv, line width=0.8pt] (proj_recv) -- ++(0, -8.6) -- (recv_query.north-|0,0);
\draw[arr=attn, line width=0.8pt] (memory) -- ++(0, -1.2) -| (recv_query);

% Key equation
\node[annotbox=3.5cm] at (9, 0.3) {
    \textbf{Focal Attention:}\\
    $h_r = \text{CrossAttn}(Q_r, K_M, V_M)$\\[2pt]
    Receiver integrates all\\
    niche information via\\
    \textbf{receiver-centered} attention
};

% ============================================
% SECTION 8: OUTPUT HEADS
% ============================================
\node[sectitle] at (-0.5, -1.2) {Output Heads};

\begin{scope}[on background layer]
    \fill[output!8, rounded corners=4pt] (-0.5, -3.0) rectangle (10.5, -1.0);
\end{scope}

% LayerNorm
\node[layernorm] (ln) at (1, -1.5) {LN};
\draw[arr=recv] (recv_query) -- (ln);

% Classification heads
\node[head] (edge1) at (3, -2.3) {Edge$_1$ Head};
\node[head] (edge2) at (5, -2.3) {Edge$_2$ Head};
\node[head] (edgen) at (7, -2.3) {$\cdots$};
\node[head, fill=output!15] (bag) at (9, -2.3) {Bag Agg};

\draw[arr=output] (ln) -- ++(0.5, 0) |- (edge1);
\draw[arr=output] (ln) -- ++(0.5, 0) |- (edge2);
\draw[arr=output!50] (ln) -- ++(0.5, 0) |- (bag);

\node[font=\tiny\sffamily, below=0.1cm of edge1] {$p(AAH\rightarrow AIS)$};
\node[font=\tiny\sffamily, below=0.1cm of edge2] {$p(AIS\rightarrow MIA)$};
\node[font=\tiny\sffamily, below=0.1cm of bag] {MIL pooling};

% ============================================
% ANNOTATIONS: Key Properties
% ============================================

% Receiver-centered highlight
\draw[recv, line width=1.5pt, rounded corners=3pt, dashed]
    (-0.3, 12.0) -- (-0.3, -0.5) -- (1.5, -0.5) -- (1.5, 10.3) -- cycle;
\node[font=\tiny\sffamily\bfseries, text=recv!80!black, rotate=90] at (-0.6, 5) {RECEIVER PATH};

% Attention flow annotation
\node[annotbox=2.5cm] at (-2, 6) {
    \textbf{Information Flow:}\\[2pt]
    1. Senders $\rightarrow$ self-attn\\
    2. L-R $\times$ Senders\\
    3. Relay self-attn\\
    4. Receiver $\leftarrow$ Memory\\[2pt]
    \textit{All flows terminate}\\
    \textit{at the receiver}
};

% Dimension key
\node[annotbox=2.2cm] at (-2, 2) {
    \textbf{Dimensions:}\\
    $d = 128$ (hidden)\\
    $k \leq 24$ (senders)\\
    $n \leq 12$ (L-R pairs)\\
    $m \leq 8$ (relay)\\
    Heads $= 4$
};

% Novel contribution box
\node[rectangle, draw=recv!70!black, fill=recv!10, rounded corners=3pt,
      font=\tiny\sffamily, align=center, text width=3cm, line width=1pt] at (9, -0.8) {
    \textbf{Novel: Attention-Weighted}\\
    \textbf{L-R Interaction Scoring}\\[2pt]
    Extract $\text{Attn}_{r\leftarrow M}$ weights\\
    to quantify biological\\
    L-R communication strength
};

% Legend
\begin{scope}[shift={(-2.3, 10.5)}]
    \node[font=\tiny\sffamily\bfseries] at (0, 0.4) {Legend:};
    \node[smalltoken=recv] at (0, 0) {};
    \node[font=\tiny\sffamily, right=0.1cm] at (0.3, 0) {Receiver};
    \node[smalltoken=sender] at (0, -0.4) {};
    \node[font=\tiny\sffamily, right=0.1cm] at (0.3, -0.4) {Sender};
    \node[smalltoken=lr] at (0, -0.8) {};
    \node[font=\tiny\sffamily, right=0.1cm] at (0.3, -0.8) {L-R Pair};
    \node[smalltoken=relay] at (0, -1.2) {};
    \node[font=\tiny\sffamily, right=0.1cm] at (0.3, -1.2) {Relay};
\end{scope}

\end{tikzpicture}
\end{document}
